\IfFileExists{stacks-project.cls}{%
\documentclass{stacks-project}
}{%
\documentclass{amsart}
}

% For dealing with references we use the comment environment
\usepackage{verbatim}
\newenvironment{reference}{\comment}{\endcomment}
%\newenvironment{reference}{}{}
\newenvironment{slogan}{\comment}{\endcomment}
\newenvironment{history}{\comment}{\endcomment}

% For commutative diagrams we use Xy-pic
\usepackage[all]{xy}

% We use 2cell for 2-commutative diagrams.
\xyoption{2cell}
\UseAllTwocells

% We use multicol for the list of chapters between chapters
\usepackage{multicol}

% This is generally recommended for better output
\usepackage{lmodern}
\usepackage[T1]{fontenc}

% For cross-file-references
\usepackage{xr-hyper}

% Package for hypertext links:
\usepackage{hyperref}

% For any local file, say "hello.tex" you want to link to please
% use \externaldocument[hello-]{hello}
\externaldocument[introduction-]{introduction}
\externaldocument[conventions-]{conventions}
\externaldocument[sets-]{sets}
\externaldocument[categories-]{categories}
\externaldocument[topology-]{topology}
\externaldocument[sheaves-]{sheaves}
\externaldocument[sites-]{sites}
\externaldocument[stacks-]{stacks}
\externaldocument[fields-]{fields}
\externaldocument[algebra-]{algebra}
\externaldocument[brauer-]{brauer}
\externaldocument[homology-]{homology}
\externaldocument[derived-]{derived}
\externaldocument[simplicial-]{simplicial}
\externaldocument[more-algebra-]{more-algebra}
\externaldocument[smoothing-]{smoothing}
\externaldocument[modules-]{modules}
\externaldocument[sites-modules-]{sites-modules}
\externaldocument[injectives-]{injectives}
\externaldocument[cohomology-]{cohomology}
\externaldocument[sites-cohomology-]{sites-cohomology}
\externaldocument[dga-]{dga}
\externaldocument[dpa-]{dpa}
\externaldocument[sdga-]{sdga}
\externaldocument[hypercovering-]{hypercovering}
\externaldocument[schemes-]{schemes}
\externaldocument[constructions-]{constructions}
\externaldocument[properties-]{properties}
\externaldocument[morphisms-]{morphisms}
\externaldocument[coherent-]{coherent}
\externaldocument[divisors-]{divisors}
\externaldocument[limits-]{limits}
\externaldocument[varieties-]{varieties}
\externaldocument[topologies-]{topologies}
\externaldocument[descent-]{descent}
\externaldocument[perfect-]{perfect}
\externaldocument[more-morphisms-]{more-morphisms}
\externaldocument[flat-]{flat}
\externaldocument[groupoids-]{groupoids}
\externaldocument[more-groupoids-]{more-groupoids}
\externaldocument[etale-]{etale}
\externaldocument[chow-]{chow}
\externaldocument[intersection-]{intersection}
\externaldocument[pic-]{pic}
\externaldocument[weil-]{weil}
\externaldocument[adequate-]{adequate}
\externaldocument[dualizing-]{dualizing}
\externaldocument[duality-]{duality}
\externaldocument[discriminant-]{discriminant}
\externaldocument[derham-]{derham}
\externaldocument[local-cohomology-]{local-cohomology}
\externaldocument[algebraization-]{algebraization}
\externaldocument[curves-]{curves}
\externaldocument[resolve-]{resolve}
\externaldocument[models-]{models}
\externaldocument[functors-]{functors}
\externaldocument[equiv-]{equiv}
\externaldocument[pione-]{pione}
\externaldocument[etale-cohomology-]{etale-cohomology}
\externaldocument[proetale-]{proetale}
\externaldocument[relative-cycles-]{relative-cycles}
\externaldocument[more-etale-]{more-etale}
\externaldocument[trace-]{trace}
\externaldocument[crystalline-]{crystalline}
\externaldocument[spaces-]{spaces}
\externaldocument[spaces-properties-]{spaces-properties}
\externaldocument[spaces-morphisms-]{spaces-morphisms}
\externaldocument[decent-spaces-]{decent-spaces}
\externaldocument[spaces-cohomology-]{spaces-cohomology}
\externaldocument[spaces-limits-]{spaces-limits}
\externaldocument[spaces-divisors-]{spaces-divisors}
\externaldocument[spaces-over-fields-]{spaces-over-fields}
\externaldocument[spaces-topologies-]{spaces-topologies}
\externaldocument[spaces-descent-]{spaces-descent}
\externaldocument[spaces-perfect-]{spaces-perfect}
\externaldocument[spaces-more-morphisms-]{spaces-more-morphisms}
\externaldocument[spaces-flat-]{spaces-flat}
\externaldocument[spaces-groupoids-]{spaces-groupoids}
\externaldocument[spaces-more-groupoids-]{spaces-more-groupoids}
\externaldocument[bootstrap-]{bootstrap}
\externaldocument[spaces-pushouts-]{spaces-pushouts}
\externaldocument[spaces-chow-]{spaces-chow}
\externaldocument[groupoids-quotients-]{groupoids-quotients}
\externaldocument[spaces-more-cohomology-]{spaces-more-cohomology}
\externaldocument[spaces-simplicial-]{spaces-simplicial}
\externaldocument[spaces-duality-]{spaces-duality}
\externaldocument[formal-spaces-]{formal-spaces}
\externaldocument[restricted-]{restricted}
\externaldocument[spaces-resolve-]{spaces-resolve}
\externaldocument[formal-defos-]{formal-defos}
\externaldocument[defos-]{defos}
\externaldocument[cotangent-]{cotangent}
\externaldocument[examples-defos-]{examples-defos}
\externaldocument[algebraic-]{algebraic}
\externaldocument[examples-stacks-]{examples-stacks}
\externaldocument[stacks-sheaves-]{stacks-sheaves}
\externaldocument[criteria-]{criteria}
\externaldocument[artin-]{artin}
\externaldocument[quot-]{quot}
\externaldocument[stacks-properties-]{stacks-properties}
\externaldocument[stacks-morphisms-]{stacks-morphisms}
\externaldocument[stacks-limits-]{stacks-limits}
\externaldocument[stacks-cohomology-]{stacks-cohomology}
\externaldocument[stacks-perfect-]{stacks-perfect}
\externaldocument[stacks-introduction-]{stacks-introduction}
\externaldocument[stacks-more-morphisms-]{stacks-more-morphisms}
\externaldocument[stacks-geometry-]{stacks-geometry}
\externaldocument[moduli-]{moduli}
\externaldocument[moduli-curves-]{moduli-curves}
\externaldocument[examples-]{examples}
\externaldocument[exercises-]{exercises}
\externaldocument[guide-]{guide}
\externaldocument[desirables-]{desirables}
\externaldocument[coding-]{coding}
\externaldocument[obsolete-]{obsolete}
\externaldocument[fdl-]{fdl}
\externaldocument[index-]{index}

% Theorem environments.
%
\theoremstyle{plain}
\newtheorem{theorem}[subsection]{Theorem}
\newtheorem{proposition}[subsection]{Proposition}
\newtheorem{lemma}[subsection]{Lemma}

\theoremstyle{definition}
\newtheorem{definition}[subsection]{Definition}
\newtheorem{example}[subsection]{Example}
\newtheorem{exercise}[subsection]{Exercise}
\newtheorem{situation}[subsection]{Situation}

\theoremstyle{remark}
\newtheorem{remark}[subsection]{Remark}
\newtheorem{remarks}[subsection]{Remarks}

\numberwithin{equation}{subsection}

% Macros
%
\def\lim{\mathop{\mathrm{lim}}\nolimits}
\def\colim{\mathop{\mathrm{colim}}\nolimits}
\def\Spec{\mathop{\mathrm{Spec}}}
\def\Hom{\mathop{\mathrm{Hom}}\nolimits}
\def\Ext{\mathop{\mathrm{Ext}}\nolimits}
\def\SheafHom{\mathop{\mathcal{H}\!\mathit{om}}\nolimits}
\def\SheafExt{\mathop{\mathcal{E}\!\mathit{xt}}\nolimits}
\def\Sch{\mathit{Sch}}
\def\Mor{\mathop{\mathrm{Mor}}\nolimits}
\def\Ob{\mathop{\mathrm{Ob}}\nolimits}
\def\Sh{\mathop{\mathit{Sh}}\nolimits}
\def\NL{\mathop{N\!L}\nolimits}
\def\CH{\mathop{\mathrm{CH}}\nolimits}
\def\proetale{{pro\text{-}\acute{e}tale}}
\def\etale{{\acute{e}tale}}
\def\QCoh{\mathit{QCoh}}
\def\Ker{\mathop{\mathrm{Ker}}}
\def\Im{\mathop{\mathrm{Im}}}
\def\Coker{\mathop{\mathrm{Coker}}}
\def\Coim{\mathop{\mathrm{Coim}}}

% Boxtimes
%
\DeclareMathSymbol{\boxtimes}{\mathbin}{AMSa}{"02}

%
% Macros for moduli stacks/spaces
%
\def\QCohstack{\mathcal{QC}\!\mathit{oh}}
\def\Cohstack{\mathcal{C}\!\mathit{oh}}
\def\Spacesstack{\mathcal{S}\!\mathit{paces}}
\def\Quotfunctor{\mathrm{Quot}}
\def\Hilbfunctor{\mathrm{Hilb}}
\def\Curvesstack{\mathcal{C}\!\mathit{urves}}
\def\Polarizedstack{\mathcal{P}\!\mathit{olarized}}
\def\Complexesstack{\mathcal{C}\!\mathit{omplexes}}
% \Pic is the operator that assigns to X its picard group, usage \Pic(X)
% \Picardstack_{X/B} denotes the Picard stack of X over B
% \Picardfunctor_{X/B} denotes the Picard functor of X over B
\def\Pic{\mathop{\mathrm{Pic}}\nolimits}
\def\Picardstack{\mathcal{P}\!\mathit{ic}}
\def\Picardfunctor{\mathrm{Pic}}
\def\Deformationcategory{\mathcal{D}\!\mathit{ef}}
\usepackage{cleveref}
\crefname{originalreference}{original reference}{original references}
\Crefname{originalreference}{Original reference}{Original references}
\makeatletter
\newlabel{example-godement-functorial}{{0G5P}{source}{}{sourcepointer-1}{}}
\newlabel{example-godement-functorial@cref}{{[originalreference][0][]0G5P}{[1][0][]source}{}{}{}}
\makeatother
\crefname{originalreference}{original reference}{original references}
\Crefname{originalreference}{Original reference}{Original references}
\makeatletter
\newlabel{lemma-godement}{{0G5N}{source}{}{sourcepointer-2}{}}
\newlabel{lemma-godement@cref}{{[originalreference][0][]0G5N}{[1][0][]source}{}{}{}}
\makeatother
\crefname{originalreference}{original reference}{original references}
\Crefname{originalreference}{Original reference}{Original references}
\makeatletter
\newlabel{lemma-relations-homotopy}{{019L}{source}{}{sourcepointer-3}{}}
\newlabel{lemma-relations-homotopy@cref}{{[originalreference][0][]019L}{[1][0][]source}{}{}{}}
\makeatother
\crefname{originalreference}{original reference}{original references}
\Crefname{originalreference}{Original reference}{Original references}
\makeatletter
\newlabel{remark-homotopy-better}{{019M}{source}{}{sourcepointer-4}{}}
\newlabel{remark-homotopy-better@cref}{{[originalreference][0][]019M}{[1][0][]source}{}{}{}}
\makeatother
\title{The two star-composites a star (b\_n) and (a\_n) star b are homotopic through a\_\{i-1\} star b\_\{n-i\}}
\author{The Stacks Project Authors}
\date{stacks-project:simplicial.tex:local-snapshot-20261012; GNU FDL 1.2 or later}
\begin{document}
\maketitle
\noindent\textit{Editorial note.} All delivered mathematical text is the author's own, copied from the pinned Stacks Project chapter file sources/stacks-simplicial/pinned-original-source.tex; no statement or proof step was rewritten, reordered or removed.\par\medskip
\noindent\textit{Editorial note.} The counted proof is the author's complete argument: the proof uses the 1\_a shorthand introduced in the proof of lemma-godement and the rewriting identities for d\textasciicircum{}n\_j and s\textasciicircum{}n\_j are displayed in-span before use; both target composites and their analogues on G' circ X circ F' are handled symmetrically.\par\medskip
\noindent\textit{Editorial note.} The label lemma-godement-two-maps names the counted statement, and the proof environment immediately following it in the source is delivered verbatim as its proof.\par\medskip
\noindent\textit{Editorial note.} The counted statement is carried by the author's own lemma environment in the chapter source.\par\medskip
\noindent\textit{Allowed standard prerequisites (editorial).}\begin{itemize}
\item The Godement functorial setup. Same chapter, \textbackslash{}label\{example-godement-functorial\}. Line 6612: the categories, Y, d, s and the composable functors F, G with X the simplicial endofunctor.
\item The relations satisfied by a simplicial homotopy. Same chapter, \textbackslash{}label\{lemma-relations-homotopy\}. Lines 6627-6632: conditions (1)-(5) that the family h\_\{n,i\} must satisfy; each is verified below.
\item The alternative homotopy formalism. Same chapter, Remark \textbackslash{}label\{remark-homotopy-better\}. Line 6631: cross-reference for the homotopy relations used.
\item The interleaving calculus for d\textasciicircum{}n\_j and s\textasciicircum{}n\_j. Same chapter, \textbackslash{}label\{lemma-godement\} (its proof introduces the 1\_a shorthand). Lines 6640-6652 and 6684-6694: d\textasciicircum{}n\_j and s\textasciicircum{}n\_j rewritten by sliding the middle d resp. s across the 1-factors.
\end{itemize}\medskip
\ifdefined\hypertarget\hypertarget{sourcepointer-1}{}\fi
\noindent\textit{Source reference.} Original reference 0G5P: sources/stacks-simplicial/pinned-original-source.tex, source lines 6563--6563. Prerequisite pointer invoked at proof line(s) 6614; it names a standard fact of the same chapter that the counted argument uses, and every step of the counted proof built on it is retained verbatim in the delivered body.\par\medskip
\ifdefined\hypertarget\hypertarget{sourcepointer-2}{}\fi
\noindent\textit{Source reference.} Original reference 0G5N: sources/stacks-simplicial/pinned-original-source.tex, source lines 6443--6443. Prerequisite pointer invoked at proof line(s) 6651; it names a standard fact of the same chapter that the counted argument uses, and every step of the counted proof built on it is retained verbatim in the delivered body.\par\medskip
\ifdefined\hypertarget\hypertarget{sourcepointer-3}{}\fi
\noindent\textit{Source reference.} Original reference 019L: sources/stacks-simplicial/pinned-original-source.tex, source lines 4704--4704. Prerequisite pointer invoked at proof line(s) 6633; it names a standard fact of the same chapter that the counted argument uses, and every step of the counted proof built on it is retained verbatim in the delivered body.\par\medskip
\ifdefined\hypertarget\hypertarget{sourcepointer-4}{}\fi
\noindent\textit{Source reference.} Original reference 019M: sources/stacks-simplicial/pinned-original-source.tex, source lines 4737--4737. Prerequisite pointer invoked at proof line(s) 6634; it names a standard fact of the same chapter that the counted argument uses, and every step of the counted proof built on it is retained verbatim in the delivered body.\par\medskip
\begin{lemma}
\label{lemma-godement-two-maps}
Let $\mathcal{A}$, $\mathcal{B}$, $\mathcal{C}$, $Y$, $d$, $s$, $F$, $G$
be as in Example \ref{example-godement-functorial}. Let
$F' : \mathcal{A} \to \mathcal{C}$ and
$G' : \mathcal{C} \to \mathcal{B}$ be two functors.
Let $(a_n) : G \circ X \to G' \circ X$ be a morphism of simplicial objects
compatible via augmentations with $a : G \to G'$.
Let $(b_n) : X \circ F \to X \circ F'$ be a morphism of simplicial objects
compatible via augmentations with $b : F \to F'$. Then the two maps
$$
a \star (b_n), (a_n) \star b : G \circ X \circ F \to G' \circ X \circ F'
$$
are homotopic.
\end{lemma}

\begin{proof}
To show the morphisms are homotopic we construct morphisms
$$
h_{n, i} : G \circ X_n \circ F \to G' \circ X_n \circ F'
$$
for $n \geq 0$ and $0 \leq i \leq n + 1$
satisfying the relations described in Lemma \ref{lemma-relations-homotopy}.
See also Remark \ref{remark-homotopy-better}. To satisfy condition (1)
of Lemma \ref{lemma-relations-homotopy} we are forced to set
$h_{n, 0} = a \star b_n$ and $h_{n , n + 1} = a_n \star b$.
Thus a logical choice is
$$
h_{n , i} = a_{i - 1} \star b_{n - i}
$$
for $1 \leq i \leq n$. Setting $a = a_{-1}$ and $b = b_{-1}$ we
see the displayed formular holds for $0 \leq i \leq n + 1$.

\medskip\noindent
Recall that
$$
d^n_j = 1_G \star 1_j \star d \star 1_{n - j} \star 1_F
$$
on $G \circ X \circ F$ where we use
the notation $1_a = 1_{Y \circ \ldots \circ Y}$ introduced
in the proof of Lemma \ref{lemma-godement}. We are going to use below
that we can rewrite this as
\begin{align*}
d^n_j
& =
d^j_j \star 1_{n - j} =
d^{j + 1}_j \star 1_{n - j} = \ldots = d^{n - 1}_j \star 1_1 \\
& =
1_j \star d^{n - j}_0 = 1_{j - 1} \star d^{n - j + 1}_1 = \ldots =
1_1 \star d^{n - 1}_{j - 1}
\end{align*}
Of course we have the analogous formulae for $d^n_j$ on $G' \circ X \circ F'$.

\medskip\noindent
We check condition (2) of Lemma \ref{lemma-relations-homotopy}.
Let $i > j$. We have to show
$$
d^n_j \circ (a_{i - 1} \star b_{n - i}) =
(a_{i - 2} \star b_{n - i}) \circ d^n_j
$$
Since $i - 1 \geq j$ we can use one of the possible descriptions of
$d^n_j$ to rewrite the left hand side as
$$
(d^{i - 1}_j \star 1_{n - i + 1}) \circ (a_{i - 1} \star b_{n - i}) =
(d^{i - 1}_j \circ a_{i - 1}) \star b_{n - i} =
(a_{i - 2} \circ d^{i - 1}_j) \star b_{n - i}
$$
Similarly the right hand side becomes
$$
(a_{i - 2} \star b_{n - i}) \circ (d^{i - 1}_j \star 1_{n - i + 1}) =
(a_{i - 2} \circ d^{i - 1}_j) \star b_{n - i}
$$
Thus we obtain the same result and (2) is checked.

\medskip\noindent
We check condition (3) of Lemma \ref{lemma-relations-homotopy}.
Let $i \leq j$. We have to show
$$
d^n_j \circ (a_{i - 1} \star b_{n - i}) =
(a_{i - 1} \star b_{n - 1 - i}) \circ d^n_j
$$
Since $j \geq i$ we may rewrite the left hand side as
$$
(1_i \star d^{n - i}_{j - i}) \circ (a_{i - 1} \star b_{n - i}) =
a_{i - 1} \star (b_{n - 1 - i} \circ d^{n - i}_{j - i})
$$
A similar manipulation shows this agrees with the right hand side.

\medskip\noindent
Recall that
$$
s^n_j = 1_G \star 1_j \star s \star 1_{n - j} \star 1_F
$$
on $G \circ X \circ F$. We are going to use below
that we can rewrite this as
\begin{align*}
s^n_j
& =
s^j_j \star 1_{n - j} = s^{j + 1}_j \star 1_{n - j - 1} = \ldots =
s^{n - 1}_j \star 1_1 \\
& =
1_j \star s^{n - j}_0 = 1_{j - 1} \star s^{n - j + 1}_1 = \ldots =
1_1 \star s^{n - 1}_{j - 1}
\end{align*}
Of course we have the analogous formulae for $s^n_j$ on $G' \circ X \circ F'$.

\medskip\noindent
We check condition (4) of Lemma \ref{lemma-relations-homotopy}.
Let $i > j$. We have to show
$$
s^n_j \circ (a_{i - 1} \star b_{n - i}) =
(a_i \star b_{n - i}) \circ s^n_j
$$
Since $i - 1 \geq j$ we can rewrite the left hand side as
$$
(s^{i - 1}_j \star 1_{n - i + 1}) \circ (a_{i - 1} \star b_{n - i}) =
(s^{i - 1}_j \circ a_{i - 1}) \star b_{n - i} =
(a_i \circ s^{i - 1}_j) \star b_{n - i}
$$
Similarly the right hand side becomes
$$
(a_i \star b_{n - i}) \circ (s^{i - 1}_j \star 1_{n - i + 1}) =
(a_i \circ s^{i - 1}_j) \star b_{n - i}
$$
as desired.

\medskip\noindent
We check condition (5) of Lemma \ref{lemma-relations-homotopy}.
Let $i \leq j$. We have to show
$$
s^n_j \circ (a_{i - 1} \star b_{n - i}) =
(a_{i - 1} \star b_{n + 1 - i}) \circ s^n_j
$$
This equality holds because both sides evaluate to
$a_{i - 1} \star (s^{n - i}_{j - i} \circ b_{n - i}) =
a_{i - 1} \star (b_{n + 1 - i} \circ s^{n - i}_{j - i})$
by exactly the same arguments as above.
\end{proof}

\end{document}
